%% =============================================================================
%% Appendix J — Human-Centered B2B+B2C Ecosystem Governance
%% Unified Orchestration with Child Dignity, Cultural Sensitivity, Consent,
%% Maturity, Role Boundaries, Fairness, and Value Realization Threads
%%
%% Consolidates supervisor's Topics 40-44 (B2C accessibility, cultural
%% sensitivity, consent/agency, longitudinal trust, unified B2B+B2C
%% orchestration) and TOP-5 strategic priorities (child dignity, ecosystem
%% maturity, role boundaries, fairness governance, value realization) into
%% one architecturally coherent ecosystem appendix.
%%
%% Scope discipline: every section ties to Governance -> Explainability ->
%% Trust -> Adoption (the locked 4-box) and to a clinician-OR-caregiver
%% beneficiary. Items failing both tests are excluded.
%% =============================================================================

\addcontentsline{toc}{section}{Appendix J Framing} % [TOC-appendix-section-insertion]
\section*{Appendix J Framing}
\addcontentsline{toc}{section}{Appendix J Framing}
\label{sec:appj_framing}

\noindent\textbf{Appendix J Scope Contract --- Conceptual Governance-Ecosystem Specification Only.}
\label{para:appj_scope_contract}
This appendix specifies a \emph{conceptual ecosystem-governance design} for caregiver-, family-, and clinician-facing AI screening support. \textbf{The appendix is NOT}: a deployed caregiver-support system; a validated therapy or reassurance intervention; a child-monitoring platform; an emotional-support service; an autonomous family-engagement tool; or a clinical-decision-support product. Where the appendix discusses caregiver trust, longitudinal engagement, family-system support, or child-dignity safeguards, these are \emph{conceptual governance requirements for future caregiver-facing systems}, not validated practice. Caregiver-facing or family-engagement applications would require separate IRB review, prospective safety validation, vulnerable-population ethical clearance, and clinician-governed oversight --- explicitly out of scope. The current dissertation's contribution is the governance-and-explainability layer; any caregiver-facing implementation pathway is Phase 2+ future work. \textbf{Prior IRB approval references} for the $n{=}131$ PLS-SEM clinician cohort and $n{=}12$ expert-panel data streams whose findings inform this conceptual ecosystem design are consolidated in Chapter~\ref{chap:methodology} Section~\ref{para:prior_irb_approval_trail} (Prior IRB Approval Trail); the Phase 2+ caregiver-facing pathway will require a separate IRB amendment as documented in Appendix~\ref{chap:appendix_indian_cohort_acquisition}.

\noindent\textbf{What this appendix is.} A consolidated ecosystem-governance specification for the B2B (provider organisation) and B2C (clinician, caregiver, child, family) sides of AI-assisted ASD-EEG screening, with the supervisor's TOP-5 cross-cutting threads --- child dignity, ecosystem maturity, role boundaries, fairness governance, and value realization --- threaded through every section.

\noindent\textbf{What this appendix is not.} It is not a sixth empirical study. It is not a procurement document. It is not a product specification. It is the human-centred ecosystem-governance layer that operationalises the four-box causal chain (Governance $\to$ Explainability $\to$ Trust $\to$ Adoption) beyond the immediate clinician dyad to include caregivers, families, children, and the organisational ecosystem in which trust must be sustained over time.

\noindent\textbf{Inclusion test (the J Rule).} Every subsection in Appendix J must satisfy two questions:
\begin{enumerate}[noitemsep,topsep=0pt]
\item Does it strengthen Governance, Explainability, Trust, or Adoption?
\item Does it identify a specific human beneficiary --- clinician, caregiver, child, or family?
\end{enumerate}
A subsection that cannot answer both is excluded. This rule is auditable: each section header is annotated with the box(es) and beneficiary it serves.

%% =========================================================================
%% J1. UNIFIED B2B+B2C ORCHESTRATION FRAMEWORK (Topic 44, keystone)
%% =========================================================================
\subsection*{J1. Unified B2B+B2C Governance Orchestration Framework}
\addcontentsline{toc}{subsection}{J1. Unified B2B+B2C Governance}

\noindent\textbf{B2C-primary screening pathway.} Before detailing the orchestration framework, Figure~\ref{fig:b2c_screening_flow} fixes the operational sequence the framework must support: caregiver-facing data capture and AI screening (steps 1--5) feeding into a B2B clinician + governance safety layer (steps 6--8). The remainder of Appendix~J operationalises each step against the eight-dimension human-centred ecosystem (J2--J10).

\input{figures/fig_b2c_screening_flow}

\label{sec:appj_orchestration}

\noindent\emph{Boxes served:} Governance, Trust, Adoption. \emph{Beneficiaries:} provider organisation (B2B), clinician + caregiver + child (B2C).

\noindent\textbf{The orchestration problem.} Most AI-assisted clinical screening systems are designed as a B2B product (sold to a hospital, configured by IT, deployed for clinicians) with caregivers and children treated as data sources or end-users-by-default. This dissertation argues that for paediatric ASD-EEG screening, the B2B and B2C layers must be \emph{jointly orchestrated}: governance controls operating at the provider level (audit, drift, registry, rollback) must be visible and consequential to the family level (caregiver consent records, child dignity safeguards, longitudinal trust signals), and vice versa.

\needspace{24\baselineskip}
\noindent Table~\ref{tab:appj_orchestration} summarises unified b2b+b2c governance orchestration for appendix review.

\paragraph[Unified B2B+B2C Governance]{Unified B2B+B2C Governance Orchestration.}



{\tablestripes\tablefontsize

\noindent Table~\ref{tab:appj_orchestration} below presents Unified B2B+B2C Governance Orchestration, laying out each row in structured form to help examiners trace the source, applicability, and downstream use at a glance. % [auto-intro-strict inserted]
\paragraph{Appj Orchestration.} % [starred-section fix]
\noindent Table~\ref{tab:appj_orchestration} below presents Unified B2B+B2C Governance Orchestration, laying out each row in structured form to help examiners trace the source, applicability, and downstream use at a glance. % [auto-intro-after-heading]



\needspace{20\baselineskip}
\begin{xltabular}{\textwidth}{@{} >{\raggedright\arraybackslash}p{3.0cm} L L @{}}
\caption[Unified B2B+B2C Governance Orchestration]{Unified B2B+B2C Governance Orchestration. Each governance control is paired with a B2B operational responsibility (the organisation must enforce it) and a B2C visible signal (the family must be able to perceive that it is in force).}\label{tab:appj_orchestration} \\
\toprule
\thead{Governance control} & \thead{B2B responsibility (organisation)} & \thead{B2C visible signal (family)} \\
\midrule
\endfirsthead
\endhead
\bottomrule
\endlastfoot
\textbf{Decision audit log}      & Per-decision immutable trail; SOC2 retention; weekly review     & Caregiver-readable summary of what the AI saw and what the clinician decided \\
\textbf{Model registry + rollback} & Version pinning; tested rollback path; deploy approvals          & Plain-language statement that the version used today is the version validated for paediatric use \\
\textbf{Drift monitoring}        & Monthly recalibration; sub-group performance dashboards            & Notification when the system has been re-validated; no silent change \\
\textbf{Bias / fairness gate}    & Pre-deploy fairness pass; quarterly fairness re-test               & Caregiver-visible note that fairness has been tested across age, sex, ethnicity \\
\textbf{HITL authority}          & Clinician override is real, logged, and never overridden by AI     & Caregiver knows the clinician, not the AI, signed the decision \\
\textbf{Consent + scope}         & Tenanted consent records; revocation honoured within 24h          & Caregiver dashboard showing what data was used, by whom, when, and how to withdraw \\
\textbf{Explainability surface}  & SHAP, citations, counterfactual stored per decision                 & Family-language explanation of the top factors (not raw SHAP) \\
\textbf{Incident + escalation}   & Runbook, on-call, post-incident review                              & Family is informed of any incident affecting their child's record \\
\end{xltabular}}
\vspace{0.4em}
\noindent\textit{Note.} The orchestration is symmetric: every B2B obligation has a B2C visible-signal partner. The rule is auditable --- if a B2B control is in place but no B2C signal exists, the ecosystem governance is incomplete.

\needspace{24\baselineskip}
\begin{figure}[!htbp]
\centering
\begin{tikzpicture}[
  node distance=1.4cm,
  layerbox/.style={draw=ggunavy, fill=ggunavy!8, rounded corners=4pt, very thick,
                   minimum width=4.6cm, minimum height=1.0cm, align=center, font=\small},
  arr/.style={-{Stealth[length=3mm]}, very thick, ggunavy}]
\node[layerbox] (b2b)  {B2B orchestration layer\\\scriptsize Provider organisation: audit, registry, drift, rollback};
\node[layerbox, below=of b2b] (mid)  {Mediating governance contracts\\\scriptsize Consent, scope, fairness gates, HITL authority};
\node[layerbox, below=of mid] (b2c)  {B2C ecosystem layer\\\scriptsize Clinician + caregiver + child + family signals};
\draw[arr] (b2b) -- (mid);
\draw[arr] (mid) -- (b2c);
\draw[arr, dashed] ([xshift=4cm]b2c.east) .. controls +(right:1.5cm) and +(right:1.5cm) .. ([xshift=4cm]b2b.east)
  node[midway, right, font=\scriptsize, align=left] {Feedback:\\caregiver concerns\\flow upward};
\end{tikzpicture}
\caption[Unified Orchestration Layers]{Unified B2B+B2C orchestration. Governance contracts in the middle layer make the provider-organisation controls (top) visible and meaningful to the family ecosystem (bottom). The dashed feedback edge ensures upward flow: caregiver concerns reach governance, not just downstream.}
\label{fig:appj_orchestration_layers}
\end{figure}

\noindent\textbf{Limitation.} The orchestration framework is conceptual; empirical validation of the feedback loop (caregiver concern $\to$ governance action) is future work. The orchestration is, however, defensible as a normative model derived from socio-technical systems theory \cite{cherns1976principles} and trust-in-automation \cite{lee2004trust}.

%% =========================================================================
%% J2. CHILD DIGNITY & CHILD-CENTRIC SAFEGUARDS (TOP-5 #1)
%% =========================================================================
\subsection*{J2. Child Dignity \& Child-Centric Safeguards}
\addcontentsline{toc}{subsection}{J2. Child Dignity \& Child-Centric Safeguards}
\label{sec:appj_child_dignity}

\noindent\emph{Boxes served:} Governance, Trust. \emph{Beneficiaries:} child, caregiver.

\noindent\textbf{The dignity premise.} A paediatric ASD-EEG screening system handles data from children who cannot themselves consent and whose developmental trajectory may be shaped by the labels the system contributes to producing. Dignity is not a soft concern; it is a governance constraint with operational consequences.

\noindent Table~\ref{tab:appj_child_dignity} summarises child dignity safeguards for appendix review.

{\tablestripes\tablefontsize
\needspace{20\baselineskip}
\begin{xltabular}{\textwidth}{@{} >{\raggedright\arraybackslash}p{3.3cm} p{6.4cm} >{\raggedright\arraybackslash}p{4.6cm} @{}}
\caption[Child Dignity Safeguards]{Child Dignity Safeguards. Each safeguard is operationalised as a governance control with a specific consequence if violated.}\label{tab:appj_child_dignity} \\
\toprule
\thead{Dignity dimension} & \thead{Operational safeguard} & \thead{Consequence if violated} \\
\midrule
\endfirsthead
\endhead
\bottomrule
\endlastfoot
\textbf{No labelling without clinician} & AI never returns ``ASD'' as a label directly to families; only a clinician assigns and communicates a diagnosis & Block: family-facing UI cannot render an AI-only label \\
\textbf{Non-permanence}          & A screening output is a moment-in-time signal, not a developmental destiny; documented in family-facing copy & Audit failure: family-facing text must include non-permanence statement \\
\textbf{Minimum-necessary data}  & Only EEG channels required for screening are processed; demographic data is access-controlled & Privacy review block before deploy \\
\textbf{No re-identification}    & EEG signal de-identification verified; no longitudinal join without explicit consent & Tenant isolation enforced (\S 41 multi-tenancy policy) \\
\textbf{Right to be forgotten}   & Caregiver-initiated deletion honoured within 30 days; audit log retained with hash only & Compliance escalation if deletion not actioned \\
\textbf{Voice (developmental-stage appropriate)} & Where age and capacity permit, child is informed in age-appropriate language about EEG procedure & Required as part of consent workflow \\
\end{xltabular}}
\vspace{0.4em}
\noindent\textit{Note.} Each safeguard maps to either a system-level enforcement (block, audit failure, tenant isolation) or a procedural enforcement (consent workflow). Dignity is governance-enforceable, not aspirational text.

\noindent\textbf{Limitation.} The dignity framework draws on the UN Convention on the Rights of the Child (Article 12: right to be heard) and existing paediatric clinical-research ethics. The dissertation does not propose new bioethical principles; it operationalises established ones as governance controls.

%% =========================================================================
%% J3. ACCESSIBILITY & INCLUSIVE CAREGIVER SUPPORT (Topic 40)
%% =========================================================================
\subsection*{J3. Accessibility \& Inclusive Caregiver Support}
\addcontentsline{toc}{subsection}{J3. Accessibility \& Inclusive Caregiver Support}
\label{sec:appj_accessibility}

\noindent\emph{Boxes served:} Explainability, Trust, Adoption. \emph{Beneficiaries:} caregiver, family.

\noindent\textbf{The accessibility premise.} A caregiver who cannot read, hear, or interact with a screening result on equal terms is excluded from the trust relationship, regardless of how well-engineered the underlying model is. Accessibility is therefore an explainability constraint, not a frontend afterthought.

\noindent Table~\ref{tab:appj_accessibility} summarises accessibility operationalisation for appendix review.

{\tablestripes\tablefontsize
\needspace{20\baselineskip}
\begin{xltabular}{\textwidth}{@{} >{\raggedright\arraybackslash}p{3.5cm} p{8.1cm} >{\centering\arraybackslash}p{2.8cm} @{}}
\caption[Accessibility Operationalisation]{Accessibility Operationalisation. Each accessibility dimension is mapped to an explainability or adoption mechanism, with a WCAG 2.1 AA / Section 508 reference where applicable.}\label{tab:appj_accessibility} \\
\toprule
\thead{Dimension} & \thead{Operationalisation} & \thead{Standard} \\
\midrule
\endfirsthead
\endhead
\bottomrule
\endlastfoot
\textbf{Visual}          & High-contrast (4.5:1+), no colour-only signalling, text resize $\geq$ 200\% & WCAG 1.4 \\
\textbf{Auditory}        & Audio explanation of screening result; captions for any video    & WCAG 1.2 \\
\textbf{Cognitive (literacy)} & Plain-language version at $\leq$ Grade 8 reading level; visual icon glossary & WCAG 3.1.5 \\
\textbf{Cognitive (numeracy)} & Confidence shown as ``low / medium / high'' not raw probability for non-numerate caregivers & --- \\
\textbf{Motor}           & Keyboard navigation; no time-out on consent screens                  & WCAG 2.1, 2.2 \\
\textbf{Language}        & Translation available for the top-5 caregiver languages in the catchment area & --- \\
\textbf{Low-bandwidth}   & Text-only fallback for caregivers on low-connectivity links          & --- \\
\textbf{Caregiver support} & In-clinic patient navigator available to walk through results        & Operational \\
\end{xltabular}}
\vspace{0.4em}
\noindent\textit{Note.} Accessibility is treated as a precondition for informed trust, not a polish task. A caregiver who cannot access the explanation cannot calibrate trust.

\noindent\textbf{Limitation.} The accessibility specification is normative; the empirical study did not test accessibility outcomes per se. Future work should pilot the family-facing UI with caregivers across the WCAG dimensions.

%% =========================================================================
%% J4. CULTURAL SENSITIVITY (Topic 41)
%% =========================================================================
\subsection*{J4. Cultural Sensitivity \& Respectful Caregiver Interaction}
\addcontentsline{toc}{subsection}{J4. Cultural Sensitivity \& Respectful Caregiver Interaction}
\label{sec:appj_culture}

\noindent\emph{Boxes served:} Trust, Adoption. \emph{Beneficiaries:} caregiver, family.

\noindent\textbf{The cultural premise.} ASD is interpreted differently across cultures: some communities frame neurodevelopmental difference as a clinical condition requiring intervention; others frame it as identity, social variation, or spiritual significance. A trust ecosystem indifferent to these frames forces a clinical vocabulary on families for whom it is alien, and trust collapses.

\needspace{24\baselineskip}
\noindent Table~\ref{tab:appj_culture} summarises cultural sensitivity operationalisation for appendix review.

\paragraph{Cultural Sensitivity Operationalisation.}




{\tablestripes\tablefontsize
\needspace{20\baselineskip}
\begin{xltabular}{\textwidth}{@{} >{\raggedright\arraybackslash}p{3.2cm} L @{}}
\caption[Cultural Sensitivity Operationalisation]{Cultural Sensitivity Operationalisation. Each cultural dimension is mapped to a specific interaction design rule grounded in cross-cultural healthcare communication literature.}\label{tab:appj_culture} \\
\toprule
\thead{Cultural dimension} & \thead{Operationalisation in clinician + caregiver workflow} \\
\midrule
\endfirsthead
\endhead
\bottomrule
\endlastfoot
\textbf{Disclosure norms}     & Some cultures expect family-collective decision-making (not individual); UI supports multi-caregiver consent flows \\
\textbf{Diagnostic vocabulary} & ``Neurodevelopmental difference'' offered alongside ``ASD'' in family-facing copy; clinician selects framing in conversation \\
\textbf{Spiritual / religious framings} & Patient navigator trained to acknowledge, not contest, religious framings; AI never adjudicates \\
\textbf{Trust in institutions} & For communities with historical reasons to distrust medical institutions, an independent advocate role is offered \\
\textbf{Language directness}  & Direct vs.\ indirect communication styles supported; clinician training material covers both \\
\textbf{Stigma sensitivity}    & Family-facing materials use neutral language; never ``risk of'' framing that activates stigma \\
\textbf{Community elders}     & In communities where elders are decision-influential, consent workflow accommodates their presence with caregiver permission \\
\end{xltabular}}
\vspace{0.4em}
\noindent\textit{Note.} Cultural sensitivity is operationalised as workflow flexibility, not as embedded AI logic. The AI does not classify families by culture; the workflow allows families to select the communication style and framing they prefer.

\noindent\textbf{Limitation.} The framework relies on clinician training and patient-navigator infrastructure, not algorithmic logic. This is deliberate --- AI should not classify cultural identity --- but it does mean cultural sensitivity is contingent on organisational investment in human roles, which is captured in the value-realization section (J10).

%% =========================================================================
%% J5. CONSENT, AGENCY, PARTICIPATORY DECISION-MAKING (Topic 42)
%% =========================================================================
\subsection*{J5. Consent \& Caregiver Agency}
\addcontentsline{toc}{subsection}{J5. Consent, Caregiver Agency \& Participatory Decision-Making}
\label{sec:appj_consent}

\noindent\emph{Boxes served:} Governance, Trust. \emph{Beneficiaries:} caregiver, child.

\noindent\textbf{The consent premise.} Consent in AI-assisted screening must extend beyond the moment of EEG capture to cover: (a) the use of the data by the AI model, (b) the use of model output in clinical decisions, (c) the retention and re-use of data for future model training, and (d) the caregiver's continuing right to revoke any of the above. A single signature at intake does not satisfy any of these.

\needspace{24\baselineskip}
\noindent Table~\ref{tab:appj_consent} summarises consent surface for appendix review.

\paragraph{Consent Surface.}




{\tablestripes\tablefontsize
\needspace{20\baselineskip}
\begin{xltabular}{\textwidth}{@{} >{\raggedright\arraybackslash}p{3.7cm} p{6.7cm} >{\centering\arraybackslash}p{3.9cm} @{}}
\caption[Consent Surface]{Consent Surface. Each consent dimension is granular, separately revocable, and audit-logged.}\label{tab:appj_consent} \\
\toprule
\thead{Consent dimension} & \thead{What the caregiver agrees to} & \thead{Revocable} \\
\midrule
\endfirsthead
\endhead
\bottomrule
\endlastfoot
\textbf{Clinical use}        & EEG captured and used for this child's screening                                & Yes (with caveat: clinician already informed) \\
\textbf{AI model inference}  & EEG run through the AI screening model; result shown to clinician               & Yes \\
\textbf{Decision support}    & AI output used as one input to clinician's decision                              & Yes (clinician can decide without AI) \\
\textbf{Future training}     & De-identified EEG used in future model training                                  & Yes, default OFF \\
\textbf{Research}            & De-identified record used in approved research                                    & Yes, default OFF \\
\textbf{Longitudinal linkage} & Data joined to future visits for the same child                                   & Yes \\
\textbf{Family access}       & Caregiver receives results via family-facing dashboard                            & Yes \\
\end{xltabular}}
\vspace{0.4em}
\noindent\textit{Note.} Default-OFF for future training and research reflects the dignity principle (J2): consent for one purpose does not extend to others. Granularity is a governance feature, not a UX inconvenience.

\noindent\textbf{Limitation.} Granular consent imposes UX complexity. The dissertation's position is that the cost is justified by the gain in caregiver agency. Future work should evaluate whether the granularity surface is comprehensible to caregivers across literacy levels.

%% =========================================================================
%% J6. LONGITUDINAL TRUST (Topic 43)
%% =========================================================================
\subsection*{J6. Longitudinal Trust \& Sustainable Caregiver Lifecycle}
\addcontentsline{toc}{subsection}{J6. Longitudinal Trust \& Sustainable Caregiver Lifecycle}
\label{sec:appj_longitudinal}

\noindent\emph{Boxes served:} Trust, Adoption. \emph{Beneficiaries:} caregiver, family.

\noindent\textbf{The longitudinal premise.} Trust is not established at the first appointment; it is cultivated across the screening lifecycle. A trust ecosystem must therefore design for repeated, low-friction touchpoints rather than a single high-stakes encounter.

\noindent Table~\ref{tab:appj_lifecycle} summarises caregiver trust lifecycle for appendix review.

{\tablestripes\tablefontsize
\needspace{20\baselineskip}
\begin{xltabular}{\textwidth}{@{} >{\raggedright\arraybackslash}p{2.6cm} L L L @{}}
\caption[Caregiver Trust Lifecycle]{Caregiver Trust Lifecycle. Five canonical phases with the governance signal, the explainability signal, and the trust-supporting interaction in each.}\label{tab:appj_lifecycle} \\
\toprule
\thead{Phase} & \thead{Governance signal} & \thead{Explainability signal} & \thead{Trust interaction} \\
\midrule
\endfirsthead
\endhead
\bottomrule
\endlastfoot
\textbf{Pre-visit}     & Plain-language description of how AI is used; consent options previewed & FAQ-level overview of factors the AI considers          & Optional pre-visit conversation with navigator \\
\textbf{In-visit}       & Granular consent captured; clinician introduces AI as decision support  & Real-time SHAP-derived explanation in family-language    & Clinician communicates result, not AI \\
\textbf{Result delivery} & Result accompanied by registered model version + fairness pass        & Top-3 contributing factors in family-language            & Patient navigator available for questions \\
\textbf{Follow-up (30--90d)} & Confirmation that no silent model change has occurred                & Caregiver-readable summary of what was found             & Channel to ask questions without re-visit \\
\textbf{Long-term}      & Annual notification of any re-validation or recall                       & Lifetime audit access via family dashboard               & Right-to-be-forgotten still in force \\
\end{xltabular}}
\vspace{0.4em}
\noindent\textit{Note.} Longitudinal trust is operationalised as a state machine, not as marketing copy. Each phase has a specific governance + explainability + interaction obligation.

\noindent\textbf{Limitation.} The lifecycle model is normative; the empirical pilot covered the in-visit and result-delivery phases only. Future studies should track the 30--90 day and long-term phases.

%% =========================================================================
%% J7. ECOSYSTEM MATURITY MODEL (TOP-5 #2)
%% =========================================================================
\subsection*{J7. Ecosystem Maturity Model}
\addcontentsline{toc}{subsection}{J7. Ecosystem Maturity Model (Human-Centered RGAIG)}
\label{sec:appj_maturity}

\noindent\emph{Boxes served:} Governance, Trust, Adoption. \emph{Beneficiaries:} all (provider, clinician, caregiver, child).

\noindent\textbf{The maturity premise.} The RGAIG Maturity Framework (Ch.~1) describes provider-organisation governance maturity in six levels. This subsection extends the same six levels with caregiver-ecosystem maturity --- because a provider can score at RGAIG Level~5 on internal governance while delivering Level~1 caregiver experience, and that asymmetry is not a defensible end-state.

\needspace{24\baselineskip}
\noindent Table~\ref{tab:appj_maturity} summarises ecosystem maturity model for appendix review.

\paragraph{Ecosystem Maturity Model.}




{\tablestripes\tablefontsize
\needspace{20\baselineskip}
\begin{xltabular}{\textwidth}{@{} >{\centering\arraybackslash}p{1.4cm} L L @{}}
\caption[Ecosystem Maturity Model]{Ecosystem Maturity Model. Each RGAIG level is paired with a caregiver-ecosystem maturity signal. A defensible system advances both axes in step.}\label{tab:appj_maturity} \\
\toprule
\thead{Level} & \thead{Provider (RGAIG) maturity} & \thead{Caregiver-ecosystem maturity (this work)} \\
\midrule
\endfirsthead
\endhead
\bottomrule
\endlastfoot
\textbf{1 Initial}        & Ad-hoc AI use; no audit; no rollback                       & No family-facing surface; result communicated verbally only \\
\textbf{2 Documented}     & Model registry; audit log exists                            & Caregiver-readable result summary available on request \\
\textbf{3 Measured}       & Drift, fairness measured; dashboards exist                  & Family dashboard exists; basic consent granularity \\
\textbf{4 Controlled}     & Pre-deploy gates enforced; HITL real                        & Granular consent + right-to-be-forgotten; cultural sensitivity workflows live \\
\textbf{5 Adaptive}       & Continuous monitoring; rapid rollback                       & Longitudinal trust touchpoints; caregiver feedback channels operational \\
\textbf{6 Operationalised} & Cross-tenant best-practice diffusion                        & Community advisory in governance; caregivers shape model retraining priorities \\
\end{xltabular}}
\vspace{0.4em}
\noindent\textit{Note.} The intent is to make asymmetry visible. A provider claiming RGAIG~5 while caregivers operate in Level~1 conditions has an audit-detectable governance gap.

\noindent\textbf{Limitation.} The caregiver-axis maturity descriptions are derived from clinical communication literature and the present dissertation's framing; large-scale validation across institutions is future work.

%% =========================================================================
%% J8. ROLE BOUNDARIES & DEPENDENCY PREVENTION (TOP-5 #3)
%% =========================================================================
\subsection*{J8. Role Boundaries \& Dependency Prevention}
\addcontentsline{toc}{subsection}{J8. Role Boundaries \& Dependency Prevention}
\label{sec:appj_boundaries}

\noindent\emph{Boxes served:} Governance, Trust. \emph{Beneficiaries:} clinician, caregiver.

\noindent\textbf{The boundaries premise.} The threat the literature documents is not that AI takes over decisions but that humans \emph{stop exercising judgement} when AI is available (automation bias; Parasuraman \& Manzey, 2010). Role boundaries are governance controls that prevent dependency formation.

\noindent Table~\ref{tab:appj_boundaries} summarises role boundaries for appendix review.

{\tablestripes\tablefontsize

\noindent Table~\ref{tab:appj_boundaries} below presents Role Boundaries, laying out each row in structured form to help examiners trace the source, applicability, and downstream use at a glance. % [auto-intro-strict inserted]
\paragraph{Appj Boundaries.} % [starred-section fix]
\noindent Table~\ref{tab:appj_boundaries} below presents Role Boundaries, laying out each row in structured form to help examiners trace the source, applicability, and downstream use at a glance. % [auto-intro-after-heading]



\needspace{20\baselineskip}
\begin{xltabular}{\textwidth}{@{} >{\raggedright\arraybackslash}p{2.4cm} L >{\raggedright\arraybackslash}p{4.0cm} @{}}
\caption[Role Boundaries]{Role Boundaries. Each role has explicit authority limits with corresponding enforcement mechanisms.}\label{tab:appj_boundaries} \\
\toprule
\thead{Role} & \thead{Authority boundary} & \thead{Enforcement} \\
\midrule
\endfirsthead
\endhead
\bottomrule
\endlastfoot
\textbf{AI model}          & Cannot label, cannot communicate to family, cannot adjust threshold itself & UI block; family-facing layer never renders an AI-only label \\
\textbf{Clinician}         & Must independently assess every case; AI is one input            & Periodic AI-disagreement audit; clinician judgement reviewed if always agreeing \\
\textbf{Caregiver}         & Always retains right to seek second opinion or refuse              & Workflow surfaces alternatives; no dark patterns \\
\textbf{Provider org}      & Cannot share data beyond consent surface; cannot retrain silently  & Audit + tenant isolation + family notification on re-validation \\
\textbf{Patient navigator} & Supports family; cannot diagnose                                   & Training + scope-of-practice document \\
\textbf{Researcher}        & Access only to consented, de-identified data                       & IRB + data-use agreement enforced \\
\end{xltabular}}
\vspace{0.4em}
\noindent\textit{Note.} Dependency prevention is operationalised as periodic clinician-AI-disagreement audit: a clinician who always agrees with the AI may be over-relying. This is monitored, not punished.

\noindent\textbf{Limitation.} Audit-of-agreement is a novel control surface; thresholds for ``too much agreement'' must be calibrated empirically across populations.

%% =========================================================================
%% J9. FAIRNESS & BIAS-AWARENESS GOVERNANCE (TOP-5 #4)
%% =========================================================================
\subsection*{J9. Fairness \& Bias-Awareness Governance}
\addcontentsline{toc}{subsection}{J9. Fairness \& Bias-Awareness Governance}
\label{sec:appj_fairness}

\noindent\emph{Boxes served:} Governance, Trust, Adoption. \emph{Beneficiaries:} all, particularly under-represented sub-populations.

\noindent\textbf{The fairness premise.} ASD presentation differs by sex (chronic under-diagnosis in girls), by ethnicity (later diagnosis in Black and Hispanic children in U.S. data), and by socio-economic status. A screening AI trained without explicit fairness gates can re-produce and amplify these disparities. Fairness is therefore a pre-deployment governance gate, not a post-hoc audit.

\needspace{24\baselineskip}
\noindent Table~\ref{tab:appj_fairness} summarises fairness gate specification for appendix review.

\paragraph{Fairness Gate Specification.}




{\tablestripes\tablefontsize
\needspace{20\baselineskip}
\begin{xltabular}{\textwidth}{@{} >{\raggedright\arraybackslash}p{3.0cm} L L @{}}
\caption[Fairness Gate Specification]{Fairness Gate Specification. Each fairness metric has a pre-deployment threshold and a continuous-monitoring threshold; failure of either triggers a defined response.}\label{tab:appj_fairness} \\
\toprule
\thead{Metric} & \thead{Pre-deploy threshold} & \thead{Production response if failed} \\
\midrule
\endfirsthead
\endhead
\bottomrule
\endlastfoot
\textbf{Disparate impact}      & $\geq$ 0.80 across sex, ethnicity, age bands & Pre-deploy block; re-train or re-balance \\
\textbf{Equal opportunity gap} & $<$ 5 percentage points TPR gap                & Pre-deploy block; investigation required \\
\textbf{Calibration parity}    & Calibration slopes within 0.10 across groups   & Re-calibration before deploy \\
\textbf{Sub-group AUC}         & No sub-group AUC $>$ 5 points below overall    & Targeted data collection required \\
\textbf{Drift fairness}        & Quarterly re-test; same thresholds              & Rollback to last-validated version; re-train \\
\textbf{Caregiver-reported}    & Family dashboard surfaces dissatisfaction signal & Investigation; not direct retraining trigger \\
\end{xltabular}}
\vspace{0.4em}
\noindent\textit{Note.} Pre-deploy is a hard gate; production thresholds trigger investigation and rollback, not silent re-training. Caregiver-reported signals are surfaced but never directly modify model behaviour, to prevent gaming.

\noindent\textbf{Limitation.} The pilot dataset ($n=131$) is too small for meaningful subgroup analysis; the fairness specification here is normative, with subgroup analysis flagged as a precondition for any scale-up beyond the pilot.

%% =========================================================================
%% J10. VALUE REALIZATION & SUSTAINABLE IMPACT (TOP-5 #5)
%% =========================================================================
\subsection*{J10. Value Realization \& Sustainable Impact}
\addcontentsline{toc}{subsection}{J10. Value Realization \& Sustainable Impact}
\label{sec:appj_value}

\noindent\emph{Boxes served:} Adoption. \emph{Beneficiaries:} all.

\noindent\textbf{The value premise.} A governance ecosystem that is well-designed but never sustained becomes shelfware. Value realization makes the governance investment defensible to the budget owner and durable across leadership change.

\noindent Table~\ref{tab:appj_value} summarises value realization categories for appendix review.

{\tablestripes\tablefontsize
\needspace{20\baselineskip}
\begin{xltabular}{\textwidth}{@{} >{\raggedright\arraybackslash}p{3.2cm} p{7.3cm} p{3.9cm} @{}}
\caption[Value Realization Categories]{Value Realization Categories. Each value category has a measurable indicator with a baseline-and-target structure familiar to provider-organisation budget owners.}\label{tab:appj_value} \\
\toprule
\thead{Value category} & \thead{Indicator} & \thead{Measurement cadence} \\
\midrule
\endfirsthead
\endhead
\bottomrule
\endlastfoot
\textbf{Clinical}         & Time to first specialist consultation; positive predictive value & Per cohort; quarterly \\
\textbf{Trust}            & Caregiver-reported trust score; clinician-reported trust score   & Quarterly survey \\
\textbf{Adoption}         & Active-clinician rate; consent-completion rate                   & Monthly dashboard \\
\textbf{Equity}           & Sub-group adoption + outcome parity                              & Quarterly \\
\textbf{Operational}      & Rollback events; incident rate; audit-finding count              & Monthly \\
\textbf{Governance}       & RGAIG level (J7 maturity); fairness-gate pass rate               & Quarterly \\
\textbf{Caregiver}        & Right-to-be-forgotten honoured within SLA; dissatisfaction signal & Monthly \\
\textbf{Sustainability}   & Cost per screening (steady-state); staff retention in screening team & Quarterly + annually \\
\end{xltabular}}
\vspace{0.4em}
\noindent\textit{Note.} Value categories are deliberately broader than clinical metrics. A system that scores well on clinical PPV but degrades caregiver trust or equity has not realised value; the dashboard makes the trade-off legible.

\noindent\textbf{Limitation.} Value realization metrics are normative; longitudinal benchmarks are absent from this pilot and must be established as the system enters future-operational-validation pathway.

%% =========================================================================
%% J. CLOSING POSITIONING
%% =========================================================================
\subsection*{Appendix J Closing Positioning}
\addcontentsline{toc}{subsection}{Appendix J Closing Positioning}

\noindent\textbf{What Appendix J adds to the dissertation.} The five Defensibility Packs in Chapters 1--5 and Appendices A--I establish the empirical and methodological core: a 4-box causal chain, six hypotheses, a piloted survey instrument, a RGAIG Maturity Framework as hero artefact, and a locked doctoral identity. Appendix J extends the ecosystem beyond the clinician dyad to include caregivers, children, and the surrounding social context --- and threads child dignity, ecosystem maturity, role boundaries, fairness governance, and value realization through all of it. Appendix J does not introduce a new research question; it operationalises the human-centred consequences of the existing one.

\noindent\textbf{What Appendix J does not do.} It does not propose a new bioethical theory. It does not extend the empirical pilot ($n=131$) to caregivers. It does not claim AI can solve disparity, cultural alienation, or dependency formation. It documents the governance contracts and workflow specifications that, if implemented, prevent the dissertation's screening system from becoming a well-engineered B2B product that fails the family ecosystem on which paediatric ASD screening depends.

\noindent\textbf{Reviewer position.} If a reviewer asks ``does this dissertation engage with the human ecosystem beyond the clinician?'' the candidate answers: ``yes --- Appendix J, ten subsections, each tied to a beneficiary and a box in the causal chain, with the supervisor's five strategic threads operationalised as governance controls.''

\subsection*{Appendix J Competitive Positioning}
\addcontentsline{toc}{subsection}{Appendix J Competitive Positioning}

\noindent\textbf{Market positioning of RGAIG.} Table~\ref{tab:competitor_comparison} positions RGAIG against five adjacent solution categories across twelve capability dimensions. The differentiation argument is the simultaneous coverage of all twelve, not any single dimension: clinical ASD screening tools have ASD evidence + HITL + EHR but no governance maturity or EEG biomarker capture; research-grade EEG has the biomarker but no caregiver-facing pathway; consumer wearables have the caregiver-facing pathway but no ASD-specific evidence or governance; telehealth ASD apps have caregiver use + ASD evidence but limited governance; generic AI governance frameworks (NIST AI RMF, ISO/IEC 42001) have the governance scaffolding but no clinical or consumer instantiation. RGAIG is the integrated solution.

\input{figures/fig_competitor_comparison}

